\documentclass[11pt]{article}
\usepackage[review]{acl}
\usepackage{times}
\usepackage{latexsym}
\usepackage{booktabs}
\usepackage{multirow}
\usepackage{graphicx}
\usepackage{amsmath}
\usepackage{amssymb}
\usepackage{xcolor}
\usepackage{url}
\usepackage{microtype}
\setlength{\textfloatsep}{8pt plus 2pt minus 2pt}
\setlength{\floatsep}{8pt plus 2pt minus 2pt}
\setlength{\intextsep}{8pt plus 2pt minus 2pt}

\newcommand{\myna}{\textsc{Myna}}
\newcommand{\stream}{\textsc{Myna-Stream}}
\newcommand{\vsix}{\textsc{Myna-v6}}
\newcommand{\ie}{i.e.,}
\newcommand{\eg}{e.g.,}

\title{\myna: Cardinality-Aware Routing, Bundle Consistency, and\\Conformal Abstention for Non-Autoregressive\\Agent Decision Engines}
\author{Anonymous ACL submission}
\date{}

\begin{document}
\maketitle

\begin{abstract}
Agent loops make many small typed decisions over states that grow. A non-autoregressive decision engine answers such questions by scoring a fixed option set instead of generating text. We study three failure modes of this design on a frozen 16-cell suite: accuracy collapse when many options are packed into one sequence, logical inconsistency across the questions of a bundle, and uncalibrated confidence. We report \myna, two engines behind one typed \texttt{choice}/\texttt{score}/\texttt{noul} contract: a 14.45M gated-linear-attention streaming engine with a 576 KiB state, and a 446.63M fine-tuned DeBERTa-v3-large with cardinality-routed readout. Routing questions with more than 16 options to decoupled scoring repairs the 77-way banking77 cell from 0.0086 to 0.5345 and lifts 16-cell macro by 0.0329, while the other 15 cells remain identical. The accuracy arm reaches 0.8583 macro (our re-score is 0.8620, 95\% CI [0.8430, 0.8800]) and 83.76\% joint bundle exact match, but it is fine-tuned in-domain while the baselines are zero-shot on 14 of 16 cells. Under an explicit entailment predicate we measure 3/172 contradictions, not the 0/172 the checkpoint records; the campaign's predicate is not released. Calibration is good (ECE 0.0189, Brier 0.0434), but 90\% conformal sets cover only 85.49\%, so we do not claim coverage. The streaming arm is 3.04--4.13$\times$ faster than a 421M-parameter cross-encoder, runs in a browser at $1.03\times10^{-6}$ probability drift, and is not the accuracy arm. The contribution is a measurement discipline: every claim carries its arm, its denominator, and what it does not license.
\end{abstract}

\section{Introduction}
\label{sec:intro}

An agent that processes tickets, threads, or page states asks many small questions: is this a refund request, how severe is it, which department owns it. Each question has a small, fixed answer set, and the state on which it is asked grows with every turn. Autoregressive language models answer such questions fluently but pay for generation on every call and can return text outside the option set. A non-autoregressive \emph{decision engine} instead returns a categorical distribution over predefined options \citep{li2026jev,duan2026gdm}. The design is attractive for high-frequency routing, but three things can go wrong, and all three are measurable.

First, \textbf{option cardinality is a budget}. A single sequence that carries a question, a state, and every option competes for one fixed window. On a 77-way intent question the option spans alone consume almost the whole window, and a packed model can fall to near chance even though it is accurate on 4-way questions. Second, \textbf{a bundle can contradict itself}: a model may call an article “sports” and then answer “no” to “is this sports?”. Third, \textbf{confidence can be miscalibrated}, and an abstention rule built on it can withhold the wrong answers.

This paper studies those three failure modes on \textsc{decision-v2-pilot}, a frozen suite of 1{,}176 bundles and 1{,}440 questions over 16 cells and 11 public sources (Table~\ref{tab:dataset}). We report \myna, which is not one model but two engines behind one typed contract. \stream{} is a 14.45M-parameter gated-linear-attention engine whose retained state is 576 KiB at every length from 128 to 16{,}384 tokens; it scans a state once and answers later questions from the cache. \vsix{} is a 446.63M-parameter fine-tuned DeBERTa-v3-large with a cardinality-routed readout; it re-encodes the state on every ask and is the accuracy arm. The two share the question schema, the answer record, the abstention seam, and the serving API, and nothing else. Keeping them separate is the first methodological decision of the paper: a speed number from \stream{} is not a statement about \vsix{}, and an accuracy number from \vsix{} is not a statement about \stream{}.

Our contributions are four measurements and one negative mechanism result.
\begin{itemize}
\item \textbf{A cardinality diagnosis and a routing fix.} We measure the token budget of the packed layout on every cell. All questions with $K\le15$ options fit with zero dropped options; the 77-way cell overflows on 50 of 116 rows, consumes 490 option-span tokens (6.36 tokens per option), and scores 0.0086. Routing $K>16$ to decoupled per-option scoring raises that cell to 0.5345 (v5) and 0.6638 (v6), and lifts 16-cell macro by 0.0329 and 0.0081 respectively. The other 15 cells are byte-identical at a macro of 0.8712; the gain is one cell, and we say so.
\item \textbf{A reproducible bundle-consistency measurement.} The checkpoint records 0/172 contradictions, but the predicate that produced it is not in the repository. We define an explicit entailment predicate, re-score the same 172 bundles, and measure 3/172 for \vsix{}, against 31/172 for Laya 0.4.19 and 29/172 for VegaML 0.15.0 on the same rows. The cross-question graph module in the checkpoint is inert at inference, so consistency is a property of the fine-tuned encoder, not an enforced constraint.
\item \textbf{A calibration and abstention audit.} Per-primitive temperature scaling gives ECE 0.0189 and Brier 0.0434 with zero argmax flips between dev and test. Selective accuracy at 60\% coverage is 98.73\% as reported (98.15\% under our re-ranking). But split-conformal 90\% sets cover 85.49\% of test questions, not 90\%, and 95\% sets cover 93.68\%; we report the undercoverage and treat the sets as a ranking signal rather than a coverage certificate.
\item \textbf{A systems study of the streaming arm.} On one Apple M5, \stream{} answers 3.04--4.13$\times$ faster than a 421.29M-parameter Laya cross-encoder end-to-end and 6.78$\times$ faster with a cached state at one question. It runs under MLX at 1.76--2.22$\times$ over PyTorch MPS and in Chrome/WebAssembly at 59.32 MiB and $1.03\times10^{-6}$ maximum probability drift. int8 is 3.17$\times$ smaller but not faster and not faithful (6.48e-2 drift in ONNX, one argmax flip in MLX), so fp32 ships.
\end{itemize}

The fifth result is negative and shapes the rest: the same architecture trained from scratch on this suite plateaus at 0.4893--0.5339, and its long-context arm is flat at chance (0.183 at 4k, 0.150 at 16k). Removing a majority-label shortcut eliminates four constant-emitter cells but costs 0.0554 macro. The gap to the fine-tuned arm is a representation gap, not a dose we can fix by training longer, and the paper does not claim otherwise.

\paragraph{What this paper is not.} It is not a state-of-the-art claim: the accuracy arm is fine-tuned in-domain, while Laya and VegaML are zero-shot on 14 of 16 cells. It is not a new backbone, and it is not a claim that \vsix{} streams: it re-encodes, and the 576 KiB state belongs to \stream{}.

\section{Related Work}
\label{sec:related}

\paragraph{Typed and zero-shot decision models.}
Cloze-style reformulation turns classification into scoring option strings \citep{schick2021pet}, and prompt-based fine-tuning with demonstrations improves few-shot performance \citep{gao2021lmbff}. Entailment-based zero-shot classification scores labels as hypotheses \citep{yin2019benchmarking}; SetFit shows that sentence-encoder fine-tuning is competitive with prompts at small label budgets \citep{tunstall2022setfit}. The recent \emph{Jev} line makes the decision-model framing explicit: a model returns a categorical distribution over predefined options without generating text \citep{li2026jev}, and JEVal benchmarks that family across 36 datasets \citep{duan2026gdm}. JEVal's findings motivate our second and third failure modes: decision models can identify the likely outcome while substantially overstating its probability, and their errors compound in multi-step systems. Our suite is smaller and our arm is in-domain, but the failure modes are the same.

\paragraph{Efficient sequence models and cascades.}
Linear attention replaces the quadratic attention matrix with a recurrence whose state is fixed in size \citep{katharopoulos2020transformers}; gated linear attention adds a data-dependent forget gate and hardware-efficient training \citep{yang2024gla}. State-space models such as Mamba \citep{gu2024mamba}, RWKV \citep{peng2023rwkv}, and RetNet \citep{sun2023retnet} make the same trade. Sparse and chunked attention \citep{beltagy2020longformer,zaheer2020bigbird,dai2019transformerxl} and IO-aware exact attention \citep{dao2022flashattention} attack the cost from other directions. Cascades and early exit decide per input how much computation to spend \citep{viola2001rapid,teerapittayanon2016branchynet,xin2020fastbert,zhou2020deebert}, and learning-to-defer routes hard cases to an expert \citep{madras2018predict}. Our two-engine design is a coarse cascade: a constant-state System-1 engine handles frequent decisions and a larger cross-encoder is the accuracy arm; we measure the two ends rather than train a joint deferral policy.

\paragraph{Calibration, selective prediction, and conformal sets.}
Temperature scaling and reliability diagrams are the standard calibration tools \citep{guo2017calibration}; modern networks remain miscalibrated under dataset shift \citep{ovadia2019uncertainty}, and deep ensembles are a strong baseline \citep{lakshminarayanan2017ensembles}. Selective classification trades coverage for risk \citep{geifman2017selective,geifman2019selectivenet}. Split conformal prediction gives distribution-free marginal coverage under exchangeability \citep{lei2018distribution,angelopoulos2021gentle}, and conformalized quantile regression reports coverage and interval length together \citep{romano2019cqr}. We follow the reporting convention, not the guarantee: our empirical coverage is below nominal, and we say so.

\paragraph{Consistency and structured readout.}
Semantic loss injects logical constraints into training \citep{xu2018semantic}; self-consistency samples multiple reasoning paths and votes \citep{wang2023selfconsistency}; guided generation constrains decoding with a finite-state machine \citep{willard2023guided}. Our setting differs: the readout is a single forward pass per question, so there is no sampling to vote over and no decoder to constrain. We measure cross-question entailment after the fact, and we report that the checkpoint's graph module does not participate in inference.

\section{The Decision Contract and Two Engines}
\label{sec:contract}

\subsection{The typed contract}
A \emph{bundle} is a state $x$ and a set of questions $Q=\{q_1,\dots,q_m\}$. Each question has a type $\tau\in\{\texttt{choice},\texttt{score},\texttt{noul}\}$ and an option set $O_q$ with cardinality $K_q=|O_q|$. A \texttt{choice} returns one option; a \texttt{score} returns a level and an expected score; a \texttt{noul} returns $p(\text{yes})$. Every answer carries a probability vector, a confidence, a 90\% conformal set, an abstention flag, and the routing mode that produced it. The contract is the same for both engines; only the scorer behind it changes.

\subsection{\stream{}: constant state}
\stream{} is a six-layer gated-linear-attention trunk with $d_{\text{model}}=384$ and 14{,}449{,}280 parameters. An observation is scanned once into a per-layer state $S\in\mathbb{R}^{H\times d_k\times d_v}$; at production width this is 589{,}824 bytes (576 KiB), independent of the observation length. \texttt{observe()} scans, \texttt{append()} scans only the new segment, and \texttt{ask()} reads the question against $S$ with a pointer head. There is no vocabulary head and no LM trunk. The engine is the one that earns the phrase “never re-reads”.

\subsection{\vsix{}: cardinality-routed cross-encoder}
\vsix{} is a fine-tuned \texttt{microsoft/deberta-v3-large} (434{,}012{,}160 parameters, 24 layers, 512-token maximum) with 12{,}619{,}022 parameters of readout, for 446{,}631{,}182 total. It keeps four heads: a packed scorer, a set-aware head for high-cardinality questions, a CORAL ordinal head that the shipped config disables, and a cross-question graph module that inference never calls. For $K\le16$ it builds one packed sequence: a task prefix, the state, then each option followed by its own \texttt{[MASK]} token; the scorer reads the hidden state at each mask. For $K>16$ it encodes each option against the state separately and scores it against the competition in its bundle. The state is truncated before the options are, because dropping context degrades one shared span while dropping an option deletes a candidate; when the option spans themselves overflow, the tail is cut, which is exactly the v5 failure we diagnose in \S\ref{sec:cardinality}. The engine re-encodes the whole state on every ask; it has no fixed-size state and no streaming property.

\subsection{What is shared, and what is not}
Both engines implement the same \texttt{observe}/\texttt{append}/\texttt{ask} surface, the same answer record, and the same fallback seam. \texttt{myna.serve} exposes \texttt{/v1/sessions} for either. The abstention rule is shared: an answer is withheld when the conformal set is empty or when the top probability is below a settable floor. The fallback records which engine committed, so a two-tier pipeline is auditable. What is not shared is the property that motivated the project: \stream{} has a constant state and \vsix{} does not. Figure~\ref{fig:arch} draws the two arms side by side; Table~\ref{tab:dataset} lists the suite.

\begin{figure*}[t]
\centering
\includegraphics[width=0.62\textwidth]{figures/fig1_architecture.png}
\caption{The two \myna{} engines behind one typed contract. \stream{} scans a growing observation once into a 576 KiB state and answers from the cache; \vsix{} re-encodes the state on every ask and routes by option cardinality. Both write the same decision record, and a low-confidence answer can be deferred to a System-2 engine with the committing engine logged. The figure is a schematic; the measurements are in \S\ref{sec:systems}.}
\label{fig:arch}
\end{figure*}

\begin{table}[t]
\centering
\footnotesize
\caption{The frozen \textsc{decision-v2-pilot} suite (11 public sources). Cells are the 16
$(\text{source},\text{question})$ pairs; $K$ is the number of options per
question. The calibration split is the one the repository commits, not the
campaign's 2{,}000-bundle fitting split (see \S\ref{sec:limitations}).}
\label{tab:dataset}
\begin{tabular}{lrrrr}
\toprule
Split & Rows & Questions & Cells & $K$ \\
\midrule
Train & 57,904 & 73,804 & 16 & 2--77 \\
Development & 1,176 & 1,440 & 16 & 2--78 \\
Calibration & 448 & 568 & 16 & 2--77 \\
Test & 1,176 & 1,440 & 16 & 2--78 \\
\bottomrule
\end{tabular}
\end{table}


\section{Cardinality Is a Budget}
\label{sec:cardinality}

\subsection{The packed layout and its failure}
The packed scorer spends one sequence on the prefix, the state, and all option spans. Each span is the option text, one \texttt{[MASK]}, and a separator. The state is truncated first, and if the option spans still exceed the window, the fallback keeps a prefix of the options and collapses the context to a single token. That fallback is not a graceful degradation: it deletes the state and the tail options at the same time.

We measured the budget on every test cell with the suite's own renderer and tokenizer (Figure~\ref{fig:cardinality}a). The 15 non-banking77 cells have $K=2$--$15$; none overflows, and every option gets a mask. The 77-way banking77 cell is different: mean option-span length is 490.3 tokens, or 6.36 tokens per option; with a median 25-token state, the total exceeds 512 on 50 of 116 rows. Across the cell the fallback keeps a mean of 75.9 options and drops up to 6, but on overflow rows it also reduces the context to one token. The v5 packed arm scores 0.0086 on the cell (1 of 80 clean rows). The collapse is a budget failure, not a tokenization artifact.

\begin{figure*}[t]
\centering
\includegraphics[width=0.82\textwidth]{figures/fig2_cardinality.pdf}
\caption{(a) The packed window is a budget: each point is a cell's median state plus option-span tokens at its maximum $K$; the 77-way cell sits on the 512-token line. (b) Routing moves exactly one cell: v5 Packed (open) to v5 Unified (filled) on the 16 cells, sorted by $K$. The other 15 cells are identical; banking77 moves from 0.0086 to 0.5345, which is $0.5259/16=0.0329$ of macro.}
\label{fig:cardinality}
\end{figure*}

\subsection{Routing}
The router sends $K\le16$ to the packed scorer and $K>16$ to the set-aware head. The threshold is a design cap, not a measured optimum: every cell at or below 15 options fits with zero drops, and the only cell above the cap is the 77-way one. We chose 16 to leave one option of headroom. On the frozen rows, routing lifts banking77 from 0.0086 to 0.5345 and the 16-cell macro from 0.8173 to 0.8502. The other 15 cells are byte-identical; their macro is 0.8712 in both arms, and the entire +0.0329 is $0.5259/16$ from banking77. The v6 set-aware head adds another +0.1293 on the cell and +0.0081 macro, again with the other 15 cells unchanged. Both steps are one-cell repairs, and Appendix Figure~\ref{fig:ablation} shows the zero deltas rather than hiding them in an aggregate.

\subsection{What routing does not buy}
Routing is not a general accuracy method. It repairs the one cell whose option spans do not fit; it cannot improve the other 15, and it does not change the model's representation. The v5-to-v6 step is also a one-cell repair: the set-aware head is only used above the cap. A reader who wants a broad accuracy gain should look to the fine-tuned encoder, and the footing paragraph before Table~\ref{tab:compact} says why that gain is not a like-for-like comparison. We report the routing result as a budget-allocation result, with its denominator visible.

\subsection{The head-to-head, with its footing}
Table~\ref{tab:compact} reports the campaign's 16-cell head-to-head. \vsix{} has the best macro, 0.8583, and the best joint bundle exact match, 0.8376; the local re-score is 0.8620 [0.8430, 0.8800], a net +5 of 1{,}440 questions, and the paired gap over Laya 0.4.19 is +0.1956 [0.1701, 0.2210]. The footing is part of the result: v5 and v6 are fine-tuned on all 16 cells, while Laya and VegaML are zero-shot on 14 of 16; only the five agnews cells and boolq are in Laya's training mixture. On those six cells the gap is 0.9264 vs 0.8937, not the 0.1956 of the full table. The remaining 0.17 is a trained-versus-untrained gap, and no column of the table licenses a like-for-like budget claim.

\begin{table*}[t]
\centering
\footnotesize
\setlength{\tabcolsep}{4pt}
\caption{The head-to-head, compact, on the same 1{,}176 bundles. \textbf{v5/v6 are
fine-tuned in-domain on all 16 cells; Laya and VegaML are zero-shot on 14 of 16}
(the five agnews cells and boolq are in their mixture). “Other 15” excludes
banking77; “six shared” is Laya's six in-mixture cells. Full table in
Appendix~\ref{app:extra}.}
\label{tab:compact}
\begin{tabular}{lrrrrr}
\toprule
 & Laya & VegaML & v5 Packed & v5 Unified & v6 \\
\midrule
16-cell macro & 0.7050 & 0.4638 & 0.8173 & 0.8502 & \textbf{0.8583} \\
overall & 0.6792 & 0.4958 & 0.8069 & 0.8493 & \textbf{0.8597} \\
joint bundle exact & 0.6182 & 0.4498 & 0.7730 & 0.8248 & \textbf{0.8376} \\
banking77 (77-way) & 0.5862 & 0.3707 & 0.0086 & 0.5345 & \textbf{0.6638} \\
other 15 macro & 0.7129 & 0.4700 & \textbf{0.8712} & \textbf{0.8712} & \textbf{0.8712} \\
six shared cells & 0.8937 & 0.4147 & \textbf{0.9264} & \textbf{0.9264} & \textbf{0.9264} \\
\bottomrule
\end{tabular}
\end{table*}


\section{Bundle Consistency}
\label{sec:consistency}

\subsection{Bundles and the entailment predicate}
A bundle can ask several questions about one state. The suite has 172 such bundles: 92 agnews bundles with a topic question and two of the four \texttt{is\_*} booleans, and 80 yelp bundles with a rating and a recommendation. The gold labels are logically consistent: an article's topic entails its own boolean and excludes the others, and a 4--5 star rating entails recommendation while a 1--2 star rating entails non-recommendation. We define an answer as a \emph{contradiction} iff the predicted topic is $c$, the bundle asks \texttt{is\_$c$}, and the predicted boolean is “no”; or the predicted rating is 4--5 and recommendation is “no”, or 1--2 and recommendation is “yes”. The predicate is deliberately narrow: it checks only the entailments the suite's own gold labels satisfy, not every possible semantic relation.

\subsection{The recorded number and the reproducible number}
The checkpoint's audit block records 0/172 contradictions on test and 1/172 on dev. We could not reproduce the test number, because the campaign's predicate is not in the repository. Under the predicate above, on the released per-row outcomes, \vsix{} has 3/172 contradictions (2 agnews, 1 yelp). The two same-row baselines are much worse: Laya 0.4.19 has 31/172 (5 agnews, 26 yelp) and VegaML 0.15.0 has 29/172 (26 agnews, 3 yelp). Joint bundle exact match is 83.93\% for \vsix{}, 56.63\% for Laya, and 51.19\% for VegaML on the same rows. The campaign's head-to-head records 83.76\% for \vsix{}, 61.82\% for Laya's typed configuration, and 44.98\% for VegaML; those columns are a different set of runs, and we keep them separate.

The discrepancy between 0/172 and 3/172 is itself a result. A number that cannot be reproduced from the released artifacts is not evidence of zero contradictions, even when it is in the model card. We report the reproducible number and the recorded one side by side, and we do not use either as a training objective claim.

\subsection{The mechanism is the encoder, not the graph}
The checkpoint carries a cross-question graph module, but inference never calls it: \texttt{\_ask} scores each question independently, and the module's gate is $\sigma(-2.3757)=0.0850$, trained down from its initialization rather than up. The CORAL ordinal head is disabled by the shipped config, so the \texttt{score} primitive is a packed argmax over the five level options. There is no consistency loss in the checkpoint config and no training code in the repository, so we make no claim about one. What we can say is narrower and more useful: the measured consistency is a property of the shared fine-tuned encoder and the state, not of an explicit constraint. The graph module is present but inert, and that is a negative mechanism result.

\section{Calibration and Conformal Abstention}
\label{sec:calibration}

\subsection{Temperature scaling}
Temperatures are fit per primitive on a held-out calibration split: 0.9 for \texttt{choice}, 0.9 for \texttt{score}, 0.85 for \texttt{noul}. They produce zero argmax flips between dev and test, so the calibration step changes probabilities but not decisions. On the 1{,}440 test questions, the reported ECE is 0.0189 and the mean Brier score is 0.0434; our 10-bin recomputation from the released outcomes gives ECE 0.0190, which matches. Figure~\ref{fig:calibration}a is the reliability diagram; the largest bin (847 questions at confidence 0.98) is accurate to within 0.3 points.

\subsection{Risk--coverage}
Selective prediction is a curve, not a point. Figure~\ref{fig:calibration}b shows four curves: all 11 calibration sources (568 questions), the G5 trio on test (348), the G5 trio on dev (348), and the G5 trio on calibration (120). At 60\% coverage the accuracies are 0.968, 0.990, 0.971, and 0.986 respectively. The gate's own trio slice reports 0.9904 (103/104) at 60\% coverage; that is a 104-question slice and its interval is wide. The reported all-question selective accuracy at 60\% is 0.9873 (853/864); re-ranking the released outcomes by top probability gives 0.9815 (848/864). We report both and do not treat the 0.58-point difference as resolved, because the campaign's selective routine is not released.

The v0 streaming arm, on the same calibration split, reaches a maximum of 0.596 at 10.0\% coverage and fails the same gate. The confidence mechanism is therefore a property of the v6 arm, not of the architecture family.

\subsection{Conformal sets and undercoverage}
The 90\% split-conformal sets are fit per primitive on the calibration split. On test, their empirical coverage is 85.49\% (1{,}231/1{,}440), and the 95\% sets cover 93.68\% (1{,}349/1{,}440). Both are below nominal. This is not a rounding error: the gap is 4.5 points at the 90\% level. Exchangeability is the assumption that licenses marginal coverage, and it is questionable here: the calibration and test splits differ in source mix, the questions inside a bundle share a state, and the option cardinality ranges from 2 to 78. We therefore report the sets as a ranking signal and an abstention seam, not as a coverage certificate. The engine's conformal set is empty on 119 of 1{,}440 test answers even with no explicit floor; those answers are ranked in every table, so the comparison is on one definition.

\begin{figure*}[t]
\centering
\includegraphics[width=0.68\textwidth]{figures/fig5_calibration.pdf}
\caption{(a) Reliability on the 1{,}440 test questions, 10 bins; marker area is bin count. (b) Risk--coverage for all-source calibration and the G5 trio on calibration, test, and dev; dotted lines mark 0.95 accuracy at 0.60 coverage. (c) Nominal versus empirical conformal coverage: 85.49\% at 90\% and 93.68\% at 95\%.}
\label{fig:calibration}
\end{figure*}

\begin{table}[t]
\centering
\footnotesize
\caption{Calibration, conformal coverage, and selective prediction. The
"reported" column is the checkpoint's audit block; "recomputed" is our
10-bin reliability calculation and top-probability ranking on the released
1{,}440 outcomes. Conformal coverage is below nominal at both levels; we report
that undercoverage rather than a coverage guarantee.}
\label{tab:calibration}
\resizebox{\columnwidth}{!}{\begin{tabular}{lrr}
\toprule
 & reported & recomputed \\
\midrule
ECE (10 bins) & 0.0189 & 0.0190 \\
mean Brier & 0.0434 & -- \\
conformal 90\% empirical & 0.8549 & -- \\
conformal 95\% empirical & 0.9368 & -- \\
selective acc @ 60\% & 0.9873 & 0.9815 \\
coverage @ $\le$5\% risk & 0.7861 & -- \\
G5 trio @ 60\% & 0.9904 & 0.9900$^{\dagger}$ \\
\midrule
\multicolumn{3}{l}{v0, same suite: max accuracy 0.596 at 10.0\% coverage} \\
\bottomrule
\end{tabular}}
\\[2pt]
{\footnotesize $^{\dagger}$209-question test trio; the 103/104 figure in the
audit block is the gate's own slice.}
\end{table}


\section{Systems: Constant State and Deployment}
\label{sec:systems}

\subsection{Latency with a constant state}
On one Apple M5, in one process, with the same state and question dict, \stream{} answers 33.08 ms at one question against Laya's 100.69 ms (3.04$\times$), 74.87 against 271.99 at five (3.63$\times$), 120.75 against 506.28 at ten (4.19$\times$), and 869.95 against 3{,}589.50 at fifty (4.13$\times$). With the state already cached, the ask-only path is 14.86 ms at one question (6.78$\times$). Ask is 10.1 ms fixed plus 9.6 ms per question with a state slope of $-10.3$ $\mu$s/token, the fixed-size state measured; observe is 360 $\mu$s/token. The ratios are the minimum of two committed runs, and the cached range is 3.65--6.78$\times$.

\subsection{MLX and the browser}
The same 14.45M checkpoint ports to MLX over identical weights, with no conversion step. MLX fp32 is 1.76--2.22$\times$ faster than PyTorch MPS from 128 to 16{,}384 state tokens; MLX int8 is 1.79--2.41$\times$ and 17.40 MiB on disk against 55.13 MiB. The ONNX export runs in Chrome 153 under WebAssembly at 59.32 MiB artifact, 73.09 MiB on the wire, a 394.4 ms cold load, and a 435.8 ms p50 for a three-decision session; six of six parity checks pass, and the maximum probability drift against the torch engine is $1.03\times10^{-6}$. The export's own parity check is $8.94\times10^{-8}$ on the decision output.

\subsection{Why int8 does not ship}
int8 is 3.17$\times$ smaller and not faster. Five of the six linear shapes this model uses are slower quantized (0.43--1.06$\times$ fp32/int8) even though int8 reads 3.6$\times$ fewer weight bytes, because the quantized pass gets 10--37 GB/s of weight read against fp32's 55--257 GB/s on the same box. It is also not faithful: on 20 real questions the MLX int8 port moves option probabilities by 9.24e-3 and flips one argmax on a question whose top two options were 0.0008 apart; the ONNX int8 artifact drifts 6.48e-2 and fails five browser checks. fp32 ships on both paths, and the int8 artifacts stay in the tree for byte-constrained uses, not for speed.

\subsection{Two engines, one API, no shared speed claim}
The systems numbers are \stream{}'s, not \vsix{}'s. \vsix{} has no ONNX export and no browser build, and it re-encodes the state, so a latency number for it would be a different experiment. The shared API is what makes the split usable: an application can run \stream{} for high-frequency decisions and defer to \vsix{} (or any System-2 model) on abstention, with the committing engine logged. We do not claim an end-to-end cascade accuracy, because no joint deferral policy was trained.

\paragraph{The streaming arm is not the accuracy arm.}
The v0 checkpoint scores 0.9525 on the synthetic split but v1b scores 0.4893--0.5339 on the real suite; an anti-prior arm removes four constant emitters and loses 0.0554 macro; the 4k long-context arm is flat at 0.150--0.183, at the 0.167 floor. Appendix~\ref{app:extra} reports the full diagnostics, with the VOID arm and the flat-curve caveat. This is why the paper keeps the two engines separate.

\section{Discussion and Limitations}
\label{sec:limitations}

\paragraph{Footing.}
\vsix{} is fine-tuned on the suite's train split, so all 16 cells are in-domain; Laya and VegaML are zero-shot on 14 of 16. The +0.1956 gap over Laya 0.4.19 is trained-versus-untrained; on the six cells in Laya's mixture it is 0.9264 vs 0.8937.

\paragraph{Scope.}
1{,}176 bundles and 1{,}440 questions are one frozen draw from 11 public datasets; the intervals are sampling error on this split, not a population over checkpoints. JEVal's 11{,}257 instances are the right scale for a general claim.

\paragraph{What we could not verify.}
The campaign's contradiction predicate, its 2{,}000-bundle CAL split, the banking77 \texttt{none\_absent} decomposition, and the v6 training objective are not in the repository; Appendix~\ref{app:extra} records each discrepancy. \vsix{}'s 512-token maximum also puts the long-context question out of its reach.

\paragraph{Ethics.}
Typed decisions are used in routing, moderation, and triage. Abstention is not a safety guarantee: the conformal sets undercover, the suite is English and public, and the labels carry their sources' biases. A deployment should log the committing engine and keep a human in the loop for high-stakes cells. We measured latency and bytes, not energy.

\section{Conclusion}
\label{sec:conclusion}

We studied a typed decision engine on a frozen 16-cell suite and separated three failure modes usually reported as one accuracy number. Cardinality is a budget: routing repairs the 77-way cell and only that cell. Consistency is measurable (3/172 against 31 and 29 on the same rows), but the checkpoint's 0/172 is not reproducible and the graph module is inert. Calibration is good (ECE 0.0189); conformal coverage is not (85.49\% at 90\%). The streaming arm is 3--4$\times$ faster with a 576 KiB state; it is not the accuracy arm.

\bibliographystyle{acl_natbib}
\bibliography{refs}

\appendix

\section{Additional Tables and Diagnostics}
\label{app:extra}

The competitor runs are Laya \citep{laya} (0.4.19) and VegaML \citep{vegaml} (0.15.0); both are software artifacts, not papers, and their versions are pinned in the witnesses.

\begin{table*}[t]
\centering
\small
\setlength{\tabcolsep}{4.5pt}
\caption{The 16-cell head-to-head on the frozen test split (1{,}176 bundles,
1{,}440 questions). All five columns are the campaign's scored runs on the same
rows. \textbf{Myna v5/v6 are fine-tuned in-domain on all 16 cells; Laya and
VegaML are zero-shot on 14 of 16} (the five agnews cells and boolq are in their
training mixture). The comparison is therefore trained-versus-untrained on most
cells, not a like-for-like budget comparison. Best per row in bold.}
\label{tab:main}
\begin{tabular}{lrrrrr}
\toprule
Cell & Laya & VegaML & v5 Packed & v5 Unified & v6 \\
\midrule
agnews/is\_business & 0.8837 & 0.3023 & \textbf{0.9535} & \textbf{0.9535} & \textbf{0.9535} \\
agnews/is\_scitech & \textbf{0.9184} & 0.3469 & \textbf{0.9184} & \textbf{0.9184} & \textbf{0.9184} \\
agnews/is\_sports & \textbf{0.9787} & 0.1915 & \textbf{0.9787} & \textbf{0.9787} & \textbf{0.9787} \\
agnews/is\_world & 0.8889 & 0.2222 & \textbf{0.9778} & \textbf{0.9778} & \textbf{0.9778} \\
agnews/topic & \textbf{0.9052} & 0.8879 & \textbf{0.9052} & \textbf{0.9052} & \textbf{0.9052} \\
amazon/stars & 0.3250 & 0.5250 & \textbf{0.6375} & \textbf{0.6375} & \textbf{0.6375} \\
banking77/intent & 0.5862 & 0.3707 & 0.0086 & 0.5345 & \textbf{0.6638} \\
boolq/answer & 0.7875 & 0.5375 & \textbf{0.8250} & \textbf{0.8250} & \textbf{0.8250} \\
contrastive/decision & 0.5969 & 0.4490 & \textbf{0.9847} & \textbf{0.9847} & \textbf{0.9847} \\
dbpedia14/category & 0.7328 & 0.7759 & \textbf{0.9569} & \textbf{0.9569} & \textbf{0.9569} \\
imdb/positive & 0.4625 & 0.5375 & \textbf{0.9500} & \textbf{0.9500} & \textbf{0.9500} \\
mnli/relation & 0.5948 & 0.5603 & \textbf{0.8362} & \textbf{0.8362} & \textbf{0.8362} \\
sst5/sentiment & 0.4875 & 0.4250 & \textbf{0.6000} & \textbf{0.6000} & \textbf{0.6000} \\
trec/answer\_type & 0.8190 & 0.3017 & \textbf{0.9569} & \textbf{0.9569} & \textbf{0.9569} \\
yelp/rating & 0.5250 & 0.5250 & \textbf{0.6625} & \textbf{0.6625} & \textbf{0.6625} \\
yelp/recommend & 0.7875 & 0.4625 & \textbf{0.9250} & \textbf{0.9250} & \textbf{0.9250} \\
\midrule
16-cell macro & 0.7050 & 0.4638 & 0.8173 & 0.8502 & \textbf{0.8583} \\
overall accuracy & 0.6792 & 0.4958 & 0.8069 & 0.8493 & \textbf{0.8597} \\
joint bundle exact & 0.6182 & 0.4498 & 0.7730 & 0.8248 & \textbf{0.8376} \\
\midrule
primitive: choice & 0.6882 & 0.5517 & 0.7773 & 0.8649 & \textbf{0.8865} \\
primitive: score & 0.4458 & 0.4917 & \textbf{0.6333} & \textbf{0.6333} & \textbf{0.6333} \\
primitive: noul & 0.7778 & 0.4206 & \textbf{0.9306} & \textbf{0.9306} & \textbf{0.9306} \\
\bottomrule
\end{tabular}
\end{table*}


\begin{figure*}[t]
\centering
\includegraphics[width=0.88\textwidth]{figures/fig4_consistency.pdf}
\caption{(a) Joint bundle exact match for the campaign's seven scored runs. (b) Contradictions on the 172 multi-question bundles under the explicit predicate of \S\ref{sec:consistency}, for the three runs with released per-row outcomes. The checkpoint records 0/172 for v6, but its predicate is unreleased; our recomputation is 3/172.}
\label{fig:consistency}
\end{figure*}

\begin{figure*}[t]
\centering
\includegraphics[width=0.9\textwidth]{figures/fig6_systems.pdf}
\caption{(a) p50 latency against the number of questions per request, same box and inputs. (b) Speedup over Laya; the cached-state path is the ask-only path. (c) MLX fp32 and int8 against PyTorch MPS over state lengths; int8 is smaller but not faster and not faithful.}
\label{fig:systems}
\end{figure*}

\begin{figure*}[t]
\centering
\includegraphics[width=0.82\textwidth]{figures/fig3_ablation.pdf}
\caption{(a) The two routing steps: v5 Packed 0.8173, v5 Unified 0.8502, v6 0.8583, and the local re-score 0.8620. (b) Per-cell accuracy change for each step. Every non-banking77 cell is 0.0000 in both steps; the aggregate gains are the 77-way cell divided by 16. The local re-score differs from the published column by +5 of 1{,}440 questions and is shown for transparency, not as a separate arm.}
\label{fig:ablation}
\end{figure*}

\begin{table*}[t]
\centering
\small
\setlength{\tabcolsep}{2.5pt}
\caption{The routing ablation. The macro column is the 16-cell unweighted
mean; "other 15" is the macro over every cell except banking77. The last three
rows are the token budget measured on the suite: option spans alone, tokens per
option, and rows whose packed sequence overflows 512 tokens.}
\label{tab:routing}
\begin{tabular}{lrrrr}
\toprule
Arm & macro & $\Delta$ & banking77 & other 15 \\
\midrule
v5 Packed & 0.8173 & -- & 0.0086 & 0.8712 \\
v5 Unified & 0.8502 & +0.0329 & 0.5345 & 0.8712 \\
v6 & 0.8583 & +0.0081 & 0.6638 & 0.8712 \\
v6 local & 0.8620 & +0.0037 & 0.6724 & 0.8747 \\
\midrule
$K$ (banking77) & \multicolumn{4}{l}{77--78 options} \\
 & \multicolumn{4}{l}{490.3 option-span tokens; 6.36 tokens/option} \\
overflow & \multicolumn{4}{l}{50/116 rows; context collapses to 1 token} \\
 & \multicolumn{4}{l}{max 6 options dropped} \\
\bottomrule
\end{tabular}
\end{table*}

\begin{table*}[t]
\centering
\small
\caption{Bundle consistency on the 172 multi-question bundles (92 agnews,
80 yelp). The top block is the campaign's head-to-head, whose contradiction
predicate is not in the repository. The bottom block is our recomputation on
the same rows with the explicit predicate of \S\ref{sec:consistency}; the three
runs in that block are the only ones with released per-row outcomes.}
\label{tab:consistency}
\begin{tabular}{llrrr}
\toprule
 & Run & exact & contra. & macro \\
\midrule
\multirow{3}{*}{campaign} & v6 & 0.8376 & 0/172 & 0.8583 \\
 & Laya (typed) & 0.6182 & 0/172 & 0.7050 \\
 & VegaML & 0.4498 & 0/172 & 0.4638 \\
\midrule
\multirow{3}{*}{same rows} & v6 & 0.8393 & 3/172 & 0.8620 \\
 & Laya 0.4.19 & 0.5663 & 31/172 & 0.6664 \\
 & VegaML 0.15.0 & 0.5119 & 29/172 & 0.6160 \\
\bottomrule
\end{tabular}
\end{table*}

\begin{table*}[t]
\centering
\small
\caption{The small-engine diagnostics. Synthetic and real columns are
different suites: 0.9525 is the synthetic test split; 0.4893--0.5339 are the
real 16-cell suite. The anti-prior \emph{off} arm stopped at step 3{,}314 of
3{,}600 and its own harness prints VOID; the numbers are reported with that
sentence attached.}
\label{tab:diagnostics}
\begin{tabular}{llrr}
\toprule
Arm & Data & macro & note \\
\midrule
v0 & synthetic & 0.9525 & G6 gate 0.94 \\
v1b & synthetic & 0.3423 & below chance floor \\
v1b & real suite & 0.4893 & majority floor 0.4331 \\
anti-prior off & real suite & 0.5339 & 4 emitters; VOID \\
anti-prior on & real suite & 0.4785 & 0 emitters \\
v6 & real suite & 0.8583 & in-domain \\
\midrule
lc1-4k, 4k state & needle & 0.1833 & readable:false \\
lc1-4k, 16k state & needle & 0.1500 & readable:false \\
\bottomrule
\end{tabular}
\end{table*}

\begin{table}[t]
\centering
\footnotesize
\setlength{\tabcolsep}{2pt}
\caption{Systems measurements on one Apple M5. Latency is p50, 10 reps,
3 warm-ups, same process, same inputs; Laya is \texttt{Agent.system\_one} at
421.29M parameters. The MLX and browser rows are the 14.45M streaming arm,
not v6.}
\label{tab:systems}
\resizebox{\columnwidth}{!}{\begin{tabular}{lrrr}
\toprule
 & Myna-Stream & Laya & ratio \\
\midrule
parameters & 14{,}449{,}280 & 421{,}293{,}827 & 0.034$\times$ \\
retained state & 576 KiB & none (re-encode) & -- \\
\midrule
1 question & 33.08 ms & 100.69 ms & 3.04$\times$ \\
5 questions & 74.87 ms & 271.99 ms & 3.63$\times$ \\
10 questions & 120.75 ms & 506.28 ms & 4.19$\times$ \\
50 questions & 869.95 ms & 3{,}589.50 ms & 4.13$\times$ \\
ask-only, 1 q & 14.86 ms & 100.69 ms & 6.78$\times$ \\
\midrule
MLX fp32 vs MPS & 1.76--2.22$\times$ & -- & 55.13 MiB \\
MLX int8 vs MPS & 1.79--2.41$\times$ & 17.40 MiB & 9.24e-3 drift \\
ONNX browser & 435.8 ms p50 & 59.32 MiB & 1.03e-6 drift \\
ONNX int8 & 18.56 MiB & 6.48e-2 drift & RED (5 checks) \\
\bottomrule
\end{tabular}}
\end{table}


\paragraph{Discrepancies with the campaign record.}
The campaign records 0/172 contradictions; under the explicit predicate of Appendix~\ref{app:predicate} we measure 3/172, and Table~\ref{tab:consistency} keeps the two blocks apart. The campaign reports a 2{,}000-bundle CAL split; the committed calibration split has 448 rows and 568 questions, and the calibration quantiles are used as shipped. The banking77 \texttt{none\_absent} decomposition is 0/12 here against the campaign's 7/12, while the cell aggregate matches at 78/116 against 77/116; the open-set gate is reconstructed from tensor shapes. The v6 training code is not released.

\subsection{Small-Engine Diagnostics}
\label{app:diagnostics}

The streaming arm is a deployment result, not an accuracy result. The v0 checkpoint scores 0.9525 on the synthetic typed-decision test split (gate 0.94). The v1b checkpoint, trained on the real suite, scores 0.3423 on the same synthetic rows, below the 0.3546 majority floor; on the real suite it scores 0.4893 (floor 0.4331), and the data-path variant scores 0.5339. A dose-matched anti-prior arm removes the majority-label shortcut: constant-emitter cells go 4 to 0, blank-instruction accuracy rises 0.3341 to 0.4315 (+0.0974), blank-options 0.2669 to 0.3674, and test macro falls 0.5339 to 0.4785 ($-0.0554$). The \emph{off} arm stopped at step 3{,}314 of 3{,}600 and its harness prints VOID; the \emph{on} arm is a continuation. The 4k truncated-backprop arm scores 0.1833 at 4{,}096 state tokens and 0.150 at 16{,}384; the 128-token rung is also 0.150 against a 0.167 uniform floor, so the harness prints \texttt{readable: false}. Steps 799 and 1199 agree exactly on all three shared rungs (9, 13, and 11 of 60), so the limiter is the mechanism, not the dose.

\begin{figure*}[t]
\centering
\includegraphics[width=0.9\textwidth]{figures/fig7_diagnostics.pdf}
\caption{(a) Synthetic versus real-suite macro for the small engine. The anti-prior arm removes four constant emitters and loses 0.0554 macro; the off arm is VOID. (b) Needle accuracy against state length for the 4k-trained arm: flat at the 0.167 floor, with the G4 thresholds (0.90 at 4k, 0.85 at 16k) for scale.}
\label{fig:diagnostics}
\end{figure*}

\section{Reproduction and Artifact Map}
\label{app:repro}

Every number in this paper is read from a committed JSON artifact or recomputed by a script in the repository. The main commands are:
\begin{itemize}
\item \texttt{python -m bench.v6\_audit --ckpt "myna-v5-prototype 2" --split test --every 1} re-scores all 1{,}176 bundles and writes the row outcomes; our run took 806 s on CPU and reproduced all 1{,}440 question outcomes bit-for-bit.
\item \texttt{python bench/significance.py --myna runs/myna\_v6\_test.rowoutcomes.json --laya runs/laya\_0419\_test.rowoutcomes.json --report runs/v6\_test.report.json --B 20000 --seed 20261031} writes the row-clustered intervals.
\item \texttt{python -m bench.risk\_coverage --ckpt "myna-v5-prototype 2"} \\ \texttt{--split test --source banking77,dbpedia14,trec} writes the trio curve.
\item \texttt{python paper/make\_figures.py} regenerates Figures~\ref{fig:cardinality}--\ref{fig:diagnostics} from the committed JSON; \texttt{python paper/make\_tables.py} regenerates every table.
\end{itemize}
The witnesses are \texttt{runs/witness-v6-master-paper.json}, \texttt{runs/witness-competitor-head-to-head-1440.json}, \texttt{runs/witness-v5-master-paper.json}, \texttt{runs/witness-hybrid-v3.json}, \texttt{runs/v6\_test.report.json}, \texttt{runs/significance\_v6.json}, \texttt{runs/risk\_coverage\_v6*.json}, \texttt{runs/latency\_matched.json}, \texttt{runs/bench\_mlx\_int8.json}, \texttt{runs/browser\_g3.json}, and the diagnostics under \texttt{runs/g6\_*} and \texttt{runs/needle\_*}. The per-cell table with intervals is Table~\ref{tab:appendix_cells}.

\section{Decision Sequence}
\label{app:sequence}

Figure~\ref{fig:sequence} gives the call sequence for one bundle. The readout is per question; consistency is measured after the fact, not enforced by the sequence.

\begin{figure*}[t]
\centering
\includegraphics[width=0.55\textwidth]{figures/figA1_sequence.png}
\caption{The \myna{} decision sequence for one bundle. The router chooses packed scoring for $K\le16$ and decoupled scoring for $K>16$; the conformal gate attaches a 90\% set and an abstention flag.}
\label{fig:sequence}
\end{figure*}

\begin{table*}[t]
\centering
\small
\setlength{\tabcolsep}{4pt}
\caption{Per-cell local re-score with row-clustered 95\% intervals
(20{,}000 replicates, seed 20261031). The interval is on the cell accuracy;
the last two columns are Laya 0.4.19's accuracy and the paired McNemar count
(Myna-only / Laya-only).}
\label{tab:appendix_cells}
\begin{tabular}{lrrrrr}
\toprule
Cell & $n$ & Myna & 95\% CI & Laya & McNemar \\
\midrule
agnews/is\_business & 43 & 0.9535 & [0.8810, 1.0000] & 0.8837 & 4:1 \\
agnews/is\_scitech & 49 & 0.9184 & [0.8333, 0.9818] & 0.9388 & 1:2 \\
agnews/is\_sports & 47 & 0.9787 & [0.9286, 1.0000] & 0.9787 & 0:0 \\
agnews/is\_world & 45 & 1.0000 & [1.0000, 1.0000] & 0.8000 & 9:0 \\
agnews/topic & 116 & 0.9052 & [0.8475, 0.9554] & 0.9138 & 6:7 \\
amazon/stars & 80 & 0.6375 & [0.5294, 0.7407] & 0.3000 & 32:5 \\
banking77/intent & 116 & 0.6724 & [0.5872, 0.7568] & 0.2759 & 51:5 \\
boolq/answer & 80 & 0.8500 & [0.7671, 0.9250] & 0.7000 & 17:5 \\
contrastive/decision & 196 & 0.9847 & [0.9653, 1.0000] & 0.5969 & 77:1 \\
dbpedia14/category & 116 & 0.9569 & [0.9174, 0.9909] & 0.7672 & 27:5 \\
imdb/positive & 80 & 0.9625 & [0.9157, 1.0000] & 0.4625 & 43:3 \\
mnli/relation & 116 & 0.8621 & [0.7966, 0.9231] & 0.5948 & 38:7 \\
sst5/sentiment & 80 & 0.6125 & [0.5050, 0.7183] & 0.3875 & 30:12 \\
trec/answer\_type & 116 & 0.9483 & [0.9048, 0.9835] & 0.7500 & 27:4 \\
yelp/rating & 80 & 0.6250 & [0.5147, 0.7312] & 0.5250 & 18:10 \\
yelp/recommend & 80 & 0.9250 & [0.8625, 0.9762] & 0.7875 & 11:0 \\
\bottomrule
\end{tabular}
\end{table*}


\section{Contradiction Predicate and Examples}
\label{app:predicate}

The predicate is implemented in \texttt{paper/scripts/consistency\_all\_engines.py}. For an agnews bundle, let $c$ be the predicted topic label and let $b_c$ be the boolean question for $c$ if it is present. A contradiction is $b_c=\text{no}$. For a yelp bundle, let $r$ be the predicted rating in $\{1,\dots,5\}$ and $y$ the predicted recommendation. A contradiction is $(r\ge4 \land \neg y)$ or $(r\le2 \land y)$; $r=3$ is unconstrained. The gold labels satisfy the predicate on every bundle. The three v6 contradictions are \texttt{agnews/test/7326} (topic business, \texttt{is\_business}=no), \texttt{agnews/test/1106} (topic world, \texttt{is\_world}=no), and \texttt{yelp/test/17802} (4 stars, recommend=no). The predicate is narrow on purpose: it checks only the entailments the suite's gold labels satisfy, not every possible semantic relation, and it does not count a multi-topic article as a contradiction.

\end{document}
